\section{Deployment notes}\label{app:deploy}
Evidence, gaps and published cost figures behind \S\ref{sec:dis-deploy}.
None of the prices or power figures below originates from this study's
measurements.

\textbf{Temporary and roadworks signals.} Portable signals at a lane
closure have no engineered plan, so a controller needing \emph{no site
configuration} is worth more there than one matching a tuned plan. By
design, the phase model is derived from the live network at start-up
(\S\ref{sec:des-overview}); measured, the leave-one-topology-out student
matched the generalist across the mixed holdout (\S\ref{sec:res-offline}),
evidence against topology memorisation in aggregate rather than parity on
the withheld layout; and the anti-starvation floor enforces alternation
deterministically except where it records a violation
(\S\ref{sec:des-shield}), one requirement of shuttle working. Three gaps
are clear: no lane-closure layout was tested; the phase model is read
once, so mid-deployment reconfiguration needs re-initialisation; and
shuttle working's clearance-time constraint, the safety-critical one,
belongs in the shield and is not implemented.

\textbf{Emergency-vehicle priority across a corridor.} Installations
without emergency-vehicle detection give an approaching vehicle no
priority, an omission \S\ref{sec:res-ev} measures at 31.1\,s per event.
The substrate exists: junctions exchange Ed25519-signed reports over an
MQTT adapter with a real broker between publisher and subscriber.
Capability is not a deployment claim: the bus has run in simulation and
over a local broker, not a field radio under loss and latency; the
scenario is scripted, with the scope \S\ref{sec:res-ev} declares, and
corridor-scale progressive preemption is designed for but unmeasured.

\textbf{Deployment economics.} A market survey of published prices
(September 2026)\footnote{Vendor list prices, council committee papers and
published trade estimates compiled in September 2026: NVIDIA's developer
price list (July 2026), a Greater London Authority estimate for signalised
crossings, TRL's SCOOT rollout evaluation, and 1NCE and UK MVNO tariffs.
None originates from this study's measurements.} places the compute node
at well under 1\% of site cost: a Jetson-class module lists at roughly
\pounds315--400 and runs this study's 0.6B student at a community-measured
$\sim$50 tokens/s within a 7--25\,W envelope, against
\pounds80,000--130,000 for a signalised crossing. The power figure matters
more than the price, since it answers the objection that my 6\,GB desktop
GPU misrepresents a cabinet power budget, on a published figure this study
did not verify. TRL's SCOOT rollout averaged a 12.8\% delay reduction, the
bar any deployment would have to clear. Enclosure, certification,
maintenance and integration dominate hardware at these ratios, so the
whole-life case needs a costing study this report's data cannot supply.
